% !TeX program = xelatex
% ============================================================
% common/preamble.tex — общая преамбула документов ЕСПД
% (ТЗ, ПМИ, ТП, РП, РО) ОП «Программная инженерия» ФКН НИУ ВШЭ.
% Геометрия с местом под основную надпись (штамп ГОСТ 19.104-78).
% ВКР использует СВОЮ преамбулу (vkr/preamble.tex) — этот файл её
% не касается. Данные приходят из common/meta.tex (\hse-макросы).
% ============================================================

% ── Шрифты (автогенерируются из настроек проекта) ───────────
% Семейства (основной/рубленый/моноширинный) задаются в настройках
% проекта и подставляются в common/fonts.tex ПЕРЕД каждой сборкой, каждый
% с независимым фолбэком на свободный клон с кириллицей. Здесь НЕ правьте —
% меняйте шрифты в настройках проекта (а устанавливайте их в Settings → Шрифты).
% \hseDefaultMono — дефолт моношрифта документа, если он не выбран в настройках.
\newcommand{\hseDefaultMono}{Courier New}
\input{../common/fonts}
% Язык оформления по project.lang. §2.3.4.2/§2.5.3.1: англоязычная работа
% оформляется на английском (главный язык — english). Русский — по умолчанию.

\usepackage[russian,english]{babel}


% ── Геометрия страницы по ГОСТ ЕСПД (с местом под штамп) ────
\usepackage[a4paper,left=20mm,right=10mm,top=25mm,bottom=40mm,footskip=26mm,headheight=30pt,headsep=5mm]{geometry}

% ── Основные пакеты ─────────────────────────────────────────
\usepackage{indentfirst}
\setlength{\parindent}{1.25cm}
\setlength{\parskip}{6pt}
\usepackage{setspace}
\onehalfspacing

% ── Типографика ─────────────────────────────────────────────
\usepackage{microtype}
\usepackage{csquotes}

\usepackage{graphicx}
\usepackage[table]{xcolor}
% ── Трассировка требований ──────────────────────────────────
% \req{ID}{Текст} — определение требования (в ячейке «Требование»
% таблицы ТЗ: печатается только текст); \reqref{ID} — ссылка на
% требование из любого документа (печатается сам идентификатор).
% Панель «Трассировка требований» связывает определения со ссылками
% по этим макросам.
\newcommand{\req}[2]{#2}
\newcommand{\reqref}[1]{\mbox{#1}}

% Жёлтые блоки-образцы: \begin{hseExample}…\end{hseExample} заливает
% содержимое (в т.ч. списки, несколько абзацев) жёлтым фоном — знак
% «это пример, замените своим». breakable — блок переносится по страницам.
\usepackage{tcolorbox}
\tcbuselibrary{breakable}
\newtcolorbox{hseExample}{
  breakable,
  colback=yellow!25,
  colframe=yellow!25,
  boxrule=0pt, arc=0pt,
  left=4pt, right=4pt, top=4pt, bottom=4pt,
  before skip=6pt, after skip=6pt,
  parbox=false
}
\usepackage{tikz}
\usetikzlibrary{arrows.meta,fit,positioning,shapes.geometric}
\usepackage{float}
\usepackage{tabularx}
\usepackage{longtable}
\usepackage{array}
\usepackage{multirow}
\usepackage{makecell}
\usepackage{booktabs}
\usepackage{adjustbox}
\usepackage{rotating}

% ── URL и гиперссылки ───────────────────────────────────────
\usepackage{url}
\usepackage{xurl}
\urlstyle{same}
\def\UrlBreaks{\do\/\do-\do.\do=\do?\do\&\do\_}
\usepackage{hyphenat}

% ── Раскраска ячеек в таблицах сравнения ────────────────────
\newcommand{\plus}{\cellcolor{green!20}+}
\newcommand{\plusstar}{\cellcolor{green!20}+*}
\newcommand{\minus}{\cellcolor{red!20}--}
\newcommand{\planned}{\cellcolor{blue!15}\textcolor{blue}{$\pm$}}

% ── Цитирование (ГОСТ-стиль, ссылки [n] в надстрочном виде) ──
\usepackage{cite}
\makeatletter
\renewcommand\@biblabel[1]{#1.}
\renewcommand{\@cite}[2]{\textsuperscript{[#1\if@tempswa , #2\fi]}}
\makeatother

\usepackage{hyperref}
\hypersetup{
    unicode=true,
    breaklinks=true,
    bookmarksopen=true,
    bookmarksnumbered=true,
    colorlinks=true,
    linkcolor=blue,
    urlcolor=blue,
    citecolor=blue,
    pdfborder={0 0 0}
}

% ── Колонтитулы / заголовки / счётчики / списки ─────────────
\usepackage{fancyhdr}
\usepackage{lastpage}
\usepackage{titlesec}
\usepackage{chngcntr}
\usepackage{enumitem}

% ── Данные проекта (автогенерируемые макросы \hse…) ─────────
% Загружаются здесь, чтобы common/commands.tex (подключается
% документом после преамбулы) мог раздать их под именами
% \projectname, \studentname и т.д.
\input{../common/meta}
